\documentclass[11pt]{article}
\usepackage{amsmath,amsthm,amssymb}
\usepackage[margin=1in]{geometry}
\usepackage[colorlinks=true,linkcolor=blue,citecolor=blue,urlcolor=blue]{hyperref}

\newtheorem{theorem}{Theorem}
\newtheorem{proposition}[theorem]{Proposition}
\newtheorem{corollary}[theorem]{Corollary}
\newtheorem{lemma}[theorem]{Lemma}
\theoremstyle{definition}
\newtheorem{definition}[theorem]{Definition}
\newtheorem{remark}[theorem]{Remark}
\newtheorem{example}[theorem]{Example}

\newcommand{\Set}{\mathrm{Set}}
\newcommand{\Poly}{\mathrm{Poly}}
\newcommand{\Comod}{\mathrm{Comod}}
\newcommand{\PRA}{\mathrm{PRA}}
\newcommand{\Cat}{\mathrm{Cat}^{\sharp}}
\newcommand{\id}{\mathrm{id}}
\newcommand{\Hom}{\mathrm{Hom}}
\newcommand{\comp}{\triangleleft}
\DeclareMathOperator{\colim}{colim}

\title{The worker/harness weld is the regular bicomodule:\\
deciding Conjecture 5.1}
\author{MacBeth\\ \small Kodamai}
\date{2026-10-17}

\begin{document}
\maketitle

\begin{center}
\fbox{\begin{minipage}{0.93\textwidth}
\small
\textbf{Author:} MacBeth (Kodamai).\quad\textbf{Date:} 2026-10-17.\\
\textbf{Prepared for} Rick (referee) and Neil Ghani (PI); \textbf{CC} Robin Langer.\\
\textbf{Content commit:} \texttt{PLACEHOLDERHASH}.\\
\textbf{Provenance.} Registry node \texttt{weld-is-regular-bicomodule} \textbf{[proved]}
in registry \texttt{weld-is-regular-bicomodule}; machine-corroboration node
\texttt{bicomodule-axioms-machine-checked} \textbf{[computed]} (script
\texttt{scratch/2026-10-16-weld-bicomodule-verify.py}). Decides Conjecture~5.1 of
\texttt{expository/2026-09-20-workers-bicomodules-pra.tex}. Source:
Spivak--Garner--Fairbanks, \emph{Polynomial Functors: A Mathematical Theory of
Interaction}, arXiv:2111.10968 (cited ``SGF'').
\end{minipage}}
\end{center}

\begin{abstract}
Let $S,T$ be sets and let $\Delta S=S\cdot y^{S}$ denote the store (costate) comonad on
$S$ --- equivalently, the indiscrete category on $S$ and a comonoid in
$(\Poly,\comp,y)$. The worker/harness weld polynomial is $B=\Delta S\comp\Delta T$ and the
worker grading is $A=\Delta S\otimes\Delta T=\Delta(S\times T)$. Conjecture~5.1 of the
workers/bicomodules note asserts three things: \textbf{(1)} $B$ carries a genuine
bicomodule structure (a horizontal morphism of $\Cat=\Comod(\Poly)$, i.e.\ a parametric
right adjoint $\Delta T\text{-}\Set\to\Delta S\text{-}\Set$); \textbf{(2)} under the SGF
slice factorization the worker $A$ is the right-adjoint \emph{content} factor; and
\textbf{(3)} the data discarded by the workers retract idempotent $e$ is the
coproduct $\Sigma$-aggregation tail. We decide all three. Clause~(1) is \textbf{TRUE}:
$B$ is the \emph{regular} $(\Delta T,\Delta S)$-bicomodule, with left coaction
$\lambda=\delta_{\Delta S}\comp\Delta T$ and right coaction $\rho=\Delta S\comp\delta_{\Delta T}$
(Theorem~\ref{thm:regular}), the five bicomodule axioms being machine-checked as
equalities of natural transformations on five $(|S|,|T|)$ shapes. Clause~(2) is
\textbf{FALSE}: computing the represented prafunctor gives $\bar F_{B}(V)=(T\times V)^{S}$,
whose SGF content factor is the \emph{power} $(-)^{S}$, not $A$ (Theorem~\ref{thm:clause2}).
Clause~(3) is \textbf{TRUE}: the $\Sigma$-aggregation tail is indexed by the branch-map
set $T^{S}=\{g\colon S\to T\}$, exactly the data $e$ forgets (Theorem~\ref{thm:clause3}).
The note's ``Gap~5.2'' (an apparent reversal of adjointness) is dissolved: $T^{S}$ is the
\emph{index} of a left-adjoint $\Sigma$, not right-adjoint exponential content. The
honest summary: the weld is a bicomodule and its retract's discarded branch data is the
SGF aggregation tail --- but state-entanglement ($A$) and data-aggregation (the
$\Sigma$-tail) are genuinely \emph{different} collapses. We do not inflate this partial
bridge into a full worker$\,=\,$content identification; clause~(2) is a refutation.
\end{abstract}

\section{Statement decided}\label{sec:statement}

Throughout, $\Poly\subset[\Set,\Set]$ is the category of polynomial functors, fully
faithful into $[\Set,\Set]$: a polynomial $p=\sum_{s\in S_{p}}y^{p[s]}$ is the functor
$p(X)=\sum_{s}X^{p[s]}$, and a morphism of polynomials \emph{is} a natural
transformation. The composition product $\comp$ is functor composition,
$(p\comp q)(X)=p(q(X))$, with unit $y$; the Dirichlet product $\otimes$ has position set
$S_{p}\times S_{q}$ and directions $p[s]\times q[t]$, unit $y$. A comonoid in
$(\Poly,\comp,y)$ is a small category (Ahman--Uustalu; SGF Prop.~5.2).

The \textbf{store comonad} $\Delta S$ (SGF Ex.~5.6) has position set $S$ and directions
$S$ at every position; as a comonad $\Delta S(W)=S\times W^{S}$, with
$\delta_{S}(s,\varphi)=(s,\,i\mapsto(i,\varphi))$ and $\varepsilon_{S}(s,\varphi)=\varphi(s)$.
As a category it is the \textbf{indiscrete groupoid} on $S$ (a unique isomorphism between
any two objects).

\begin{definition}[the weld and the worker]
Fix sets $S,T$. The \emph{weld} polynomial is $B:=\Delta S\comp\Delta T$. The \emph{worker}
grading is $A:=\Delta S\otimes\Delta T=\Delta(S\times T)$.
\end{definition}

Conjecture~5.1 of \texttt{expository/2026-09-20-workers-bicomodules-pra.tex} asserts:
\begin{enumerate}
\item[\textbf{(1)}] \emph{(existence)} $B$ carries a $(c,d)$-bicomodule structure for some
comonoids $c,d$ --- a horizontal morphism of $\Cat=\Comod(\Poly)$, i.e.\ a parametric
right adjoint (PRA) functor $c\text{-}\Set\to d\text{-}\Set$ by the SGF correspondence
(SGF Thm.~4.28 / Prop.~3.9);
\item[\textbf{(2)}] \emph{(identification)} under the SGF slice factorization (Prop.~3.6)
of the represented prafunctor, the worker $A$ is the right-adjoint \emph{content} factor
$F/1$;
\item[\textbf{(3)}] \emph{(tail)} the data discarded by the workers retract idempotent
$e$ (the non-constant branch maps $g\colon S\to T$) is the coproduct
$\Sigma$-aggregation tail.
\end{enumerate}

\begin{theorem}[verdict]\label{thm:verdict}
With $c=\Delta T$, $d=\Delta S$:
\begin{center}
\renewcommand{\arraystretch}{1.25}
\begin{tabular}{p{0.30\textwidth}p{0.12\textwidth}p{0.48\textwidth}}
\hline
\textbf{clause} & \textbf{verdict} & \textbf{content}\\
\hline
(1) existence & \textbf{TRUE} & $B$ is the regular $(\Delta T,\Delta S)$-bicomodule;
general comonoid nonsense (Lemma~\ref{thm:reg-lemma}, Theorem~\ref{thm:regular}).\\
(2) worker $=$ content & \textbf{FALSE} & the content factor is the power $(-)^{S}$, not
$A$ (Theorem~\ref{thm:clause2}).\\
(3) discarded $=$ $\Sigma$-tail & \textbf{TRUE} & the $\Sigma$-tail is $\Sigma_{g\in T^{S}}$,
indexed by the branch maps (Theorem~\ref{thm:clause3}).\\
\hline
\end{tabular}
\end{center}
In addition, the note's \emph{Gap~5.2} (``apparent reversal of adjointness'') is
\emph{dissolved} (Theorem~\ref{thm:gap}): the branch-map set $T^{S}$ is the indexing set
of a genuine left adjoint $\Sigma$, not right-adjoint exponential content.
\end{theorem}

The honest one-line summary: \emph{the weld is a bicomodule, and the retract's discarded
branch data is exactly the SGF aggregation tail --- but the worker $A$ is a different
collapse from the bicomodule content.} State-entanglement $\neq$ data-aggregation.

\section{Clause (1): the weld is the regular bicomodule --- TRUE}\label{sec:clause1}

Recall (SGF Def.~4.2, Leibniz order of Rmk.~4.3) that a \textbf{$(c,d)$-bicomodule}
$m\colon c\to d$ is a triple $(m,\lambda,\rho)$ with a left $d$-coaction
$\lambda\colon m\to d\comp m$ and a right $c$-coaction $\rho\colon m\to m\comp c$ satisfying
\begin{itemize}
\item[(34)] left counit $(\varepsilon_{d}\comp m)\circ\lambda=\id$; left coassociativity
$(\delta_{d}\comp m)\circ\lambda=(d\comp\lambda)\circ\lambda$;
\item[(35)] right counit $(m\comp\varepsilon_{c})\circ\rho=\id$; right coassociativity
$(\rho\comp c)\circ\rho=(m\comp\delta_{c})\circ\rho$;
\item[(36)] compatibility $(\lambda\comp c)\circ\rho=(d\comp\rho)\circ\lambda$.
\end{itemize}
By SGF Thm.~4.28 / Prop.~3.9, $\Comod(\Poly)(c,d)\simeq\PRA(c\text{-}\Set,d\text{-}\Set)$:
a bicomodule $m\colon c\to d$ \emph{is} a parametric right adjoint $c\text{-}\Set\to d\text{-}\Set$.

\begin{lemma}[regular bicomodule]\label{thm:reg-lemma}
For any comonoids $c=(c,\varepsilon_{c},\delta_{c})$ and $d=(d,\varepsilon_{d},\delta_{d})$
in a monoidal category $(V,\comp,y)$, the object $d\comp c$ is a $(c,d)$-bicomodule with
\[
\lambda:=\delta_{d}\comp c\colon d\comp c\to d\comp d\comp c=d\comp(d\comp c),\qquad
\rho:=d\comp\delta_{c}\colon d\comp c\to d\comp c\comp c=(d\comp c)\comp c.
\]
\end{lemma}

\begin{proof}
Axioms (34)/(35): $\lambda$ is $\delta_{d}$ whiskered on the right by $c$, so the
left-counit and left-coassociativity laws are the comonoid laws
$(\varepsilon_{d}\comp d)\circ\delta_{d}=\id$ and
$(\delta_{d}\comp d)\circ\delta_{d}=(d\comp\delta_{d})\circ\delta_{d}$ of $d$, whiskered by
$\comp c$; likewise (35) is the comonoid laws of $c$ whiskered by $d\comp$. Axiom (36):
$\lambda$ acts on the first ($d$) tensor factor and $\rho$ on the last ($c$) factor; since
$\comp$ is a bifunctor, the two whiskerings commute (interchange law), and both
$(\lambda\comp c)\circ\rho$ and $(d\comp\rho)\circ\lambda$ equal the single map
$\delta_{d}\comp\delta_{c}\colon d\comp c\to d\comp d\comp c\comp c$.
\end{proof}

Lemma~\ref{thm:reg-lemma} is standard --- the comonoid $\comp$-product of two comonoids is
the regular bicomodule between them. Its consequence here is the point.

\begin{theorem}[clause (1)]\label{thm:regular}
The weld $B=\Delta S\comp\Delta T$ carries the regular $(\Delta T,\Delta S)$-bicomodule
structure ($c=\Delta T$, $d=\Delta S$, carrier $d\comp c=\Delta S\comp\Delta T$), with
$\lambda=\delta_{\Delta S}\comp\Delta T$ and $\rho=\Delta S\comp\delta_{\Delta T}$. Hence $B$ is
a horizontal morphism of $\Cat$, i.e.\ a prafunctor $\Delta T\text{-}\Set\to\Delta S\text{-}\Set$.
\emph{Conjecture~5.1 clause (1): TRUE.}
\end{theorem}

\begin{proof}
Immediate from Lemma~\ref{thm:reg-lemma} with $c=\Delta T$, $d=\Delta S$; both $\Delta S$
and $\Delta T$ are comonoids (store comonads).
\end{proof}

\paragraph{Machine check.}
All five axioms --- (34) counit and coassociativity, (35) counit and coassociativity,
(36) compatibility --- were verified as equalities of natural transformations (Godement
calculus on nested store-tuples) for
$(|S|,|T|)\in\{(2,2),(2,3),(3,2),(1,2),(2,1)\}$ on test objects of size $2$ and $3$; every
case returned \texttt{True}. Script
\texttt{scratch/2026-10-16-weld-bicomodule-verify.py}; registry node
\texttt{bicomodule-axioms-machine-checked} \textbf{[computed]}. This is independent
corroboration of Theorem~\ref{thm:regular}, not a dependency: the theorem stands on the
proved Lemma~\ref{thm:reg-lemma}.

\paragraph{Which two comonoids?}
This dissolves the note's worry (``Gap~5.3'') that $\Delta S\comp\Delta T$ is a
\emph{monoidal product} of two comonoids rather than, \emph{a priori}, a bicomodule
\emph{between} two fixed ones. The regular-bicomodule construction turns the $\comp$-product
of \emph{any} two comonoids into a bicomodule between them; the pair is
$(c,d)=(\Delta T,\Delta S)$.

\section{The represented prafunctor}\label{sec:represented}

Because $\Delta S$ is the indiscrete groupoid on $S$, it is equivalent as a category to the
terminal category $\mathbf 1$ (for $S\neq\varnothing$: a unique isomorphism between any two
objects). Hence $\Delta S\text{-}\Set\simeq\Set$: a copresheaf on $\Delta S$ is a single set
with canonical identifications, and $\lim X\cong X(s_{0})\cong\colim X$ for every object
$s_{0}$. Write $E_{c}\colon c\text{-}\Set\simeq\Set$ for this equivalence ($E_{c}^{-1}V$ is
the constant copresheaf $V$). We transport the prafunctor of $B$ to an endofunctor
$\bar F_{B}$ of $\Set$; right-adjointness and the SGF factorization are
equivalence-invariant.

\begin{theorem}[represented prafunctor]\label{thm:represented}
$\bar F_{B}(V)\cong(T\times V)^{S}=T^{S}\times V^{S}$. As a polynomial endofunctor of
$\Set$, $\bar F_{B}=T^{S}\cdot y^{S}$ (positions $T^{S}$, directions $S$). Concretely
\[
\bar F_{B}(X)=\{(g,h)\colon g\colon S\to T,\ h\colon S\to X\}.
\]
\end{theorem}

\begin{proof}
The regular bicomodule $d\comp c$ is the composite over the unit $y$ of two one-sided
comodules:
\[
\Delta T\ \xrightarrow{\ G:=(\Delta T\to y)\ }\ y=\Set\ \xrightarrow{\ H:=(y\to\Delta S)\ }\ \Delta S,
\qquad\text{carrier } H\comp G=\Delta S\comp\Delta T=B,
\]
where $G$ is the right $\Delta T$-comodule with carrier $\Delta T$ and
$\rho=\delta_{\Delta T}$, and $H$ is the left $\Delta S$-comodule with carrier $\Delta S$ and
$\lambda=\delta_{\Delta S}$. Bicomodule composition over $y$ is just $\comp$ (no equalizer,
since $y$ is the unit), and $H\comp G=\Delta S\comp\Delta T=B$. Prafunctors compose, so
$F_{B}=F_{H}\circ F_{G}$.

\emph{Factor $F_{G}\colon\Delta T\text{-}\Set\to\Set$.} As a $\Delta T$-variable polynomial
(SGF Prop.~3.6, $\Set[c]$), the right $\Delta T$-comodule with carrier $\Delta T$ is
$F_{G}(X)=\sum_{t\in T}\Delta T\text{-}\Set(y_{t},X)=\sum_{t\in T}X(t)$, where
$y_{t}=\Delta T(t,-)$ is the representable at $t$. In the indiscrete category $y_{t}$ is the
terminal copresheaf $\mathbf 1$ (as $\Delta T(t,t')=1$), so $\Delta T\text{-}\Set(y_{t},X)=\lim X$.
Hence $F_{G}(X)=\sum_{t\in T}\lim X$, and after transport $\bar G(V)=T\times V$.

\emph{Factor $F_{H}\colon\Set\to\Delta S\text{-}\Set$.} A left $\Delta S$-comodule is a functor
$\Delta S\to\Poly$ (SGF Prop.~5.7); the regular one (carrier $\Delta S$,
$\lambda=\delta_{\Delta S}$) assigns to each object $a\in S$ the component polynomial
$(\Delta S)_{a}=y^{S}$ (one position over $a$, with $S$ directions --- read off
$\delta_{S}$). So $F_{H}(W)(a)=(y^{S})(W)=W^{S}$ (constant copresheaf), and after transport
$\bar H(W)=W^{S}$.

Therefore $\bar F_{B}(V)=\bar H(\bar G(V))=(T\times V)^{S}=T^{S}\times V^{S}$, and
$(T\times V)^{S}=\{S\to T\times V\}=\{(g,h)\colon g\colon S\to T,\ h\colon S\to V\}$.
\end{proof}

\paragraph{The point of the whole problem.}
$\bar F_{B}=T^{S}\,y^{S}$ is \emph{neither} the underlying polynomial of
$A=(S\times T)\,y^{S\times T}$ \emph{nor} that of $B=(S\cdot T^{S})\,y^{S\times T}$. The
represented prafunctor of a bicomodule is \emph{not} the underlying polynomial functor of
its carrier --- exactly the trap the note flagged. Machine-checked cardinalities at
$|S|=|T|=2$, $V=2$: $\bar F_{B}=16$, $A_{\text{under}}=64$, $B_{\text{under}}=128$.

\section{The SGF factorization, and clauses (2) and (3)}\label{sec:factorization}

Apply the SGF slice factorization (Prop.~3.6) to
$\bar F_{B}(V)=\sum_{g\in T^{S}}V^{S}$:
\begin{itemize}
\item $\bar F_{B}(\mathbf 1)=T^{S}$;
\item \textbf{content} $\bar F_{B}/\mathbf 1\colon\Set\to\Set^{T^{S}}$,
$V\mapsto(V^{S})_{g\in T^{S}}=\Delta_{T^{S}}\circ(-)^{S}$. Both the power $(-)^{S}$ and the
diagonal $\Delta_{T^{S}}$ are right adjoints, so $\bar F_{B}/\mathbf 1$ is a right adjoint;
the genuine connected-limit content is the \textbf{power} $(-)^{S}$;
\item \textbf{$\Sigma$-tail} $\Sigma\colon\Set^{T^{S}}\to\Set$, the coproduct over
$T^{S}=\{g\colon S\to T\}$.
\end{itemize}

\begin{theorem}[clause (2): FALSE]\label{thm:clause2}
The right-adjoint content factor of the weld prafunctor is the power $(-)^{S}$, not the
worker $A$. Indeed $A=\Delta(S\times T)$ has represented functor $V\mapsto V^{S\times T}$
(exponent $S\times T$), whereas the content has exponent $S$; these differ whenever
$|T|\geq2$. No handedness convention repairs this: the opposite reading gives content
$(-)^{T}$, still a power, never $A$.
\end{theorem}

\begin{proof}
The content exponent is $S$ (the directions of $\bar F_{B}=T^{S}y^{S}$); $A$'s represented
functor has exponent $S\times T$ by the computation of Theorem~\ref{thm:represented}
applied to $A=\Delta(S\times T)$, giving $\bar A(V)=V^{S\times T}$. For $|T|\geq2$ we have
$|S|<|S\times T|$, so the two functors differ on any $V$ with $|V|\geq2$.
\end{proof}

\begin{theorem}[clause (3): TRUE]\label{thm:clause3}
The SGF $\Sigma$-aggregation tail is the coproduct $\Sigma_{g\in T^{S}}$ indexed by the
full branch-map set $T^{S}=\{g\colon S\to T\}$ --- precisely the ``non-constant branch
maps'' that the retract idempotent $e$ of the workers theorem forgets.
\end{theorem}

\begin{proof}
The tail index is $\bar F_{B}(\mathbf 1)=T^{S}$ by Theorem~\ref{thm:represented};
$T^{S}=\Hom(S,T)$ is the set of branch maps, and $e(s,g)=(s,\mathrm{const}_{g(s)})$
collapses each $g$ to a constant, forgetting exactly the non-constant elements of $T^{S}$.
\end{proof}

\begin{theorem}[Gap 5.2 dissolved]\label{thm:gap}
There is no adjointness reversal. The branch maps $T^{S}$ appear as the \emph{indexing
set} of the left-adjoint coproduct $\Sigma_{T^{S}}$, not as right-adjoint exponential
content. ``$T^{S}$ is an exponential object'' is a statement about the \emph{set}, not
about the adjointness of $\Sigma_{T^{S}}\colon\Set^{T^{S}}\to\Set$, which is left adjoint to
the diagonal. The genuinely right-adjoint (exponential / connected-limit) content
$(-)^{S}$ is \emph{kept} in $F/\mathbf 1$, exactly as SGF predicts.
\end{theorem}

\begin{proof}
$\Sigma_{T^{S}}\dashv\Delta_{T^{S}}$, so $\Sigma_{T^{S}}$ is a left adjoint for any index
set, exponential or not; and $(-)^{S}$ is right adjoint to the copower $S\times(-)$, hence
connected-limit-preserving content. The note's Gap~5.2 conflated the two.
\end{proof}

\section{The precise role of the worker}\label{sec:role}

The worker $A$ is not the content, but it has an exact location inside the weld
prafunctor. The section $\sigma\colon A\to B$, $\sigma_{\mathbf 1}(s,t)=(s,\mathrm{const}_{t})$,
embeds $A$'s $T$-coordinate as the \emph{constant} branch maps $T\hookrightarrow T^{S}$. In
the factorization of $\bar F_{B}$, the current input state $S$ lives in the content exponent
and the branch map lives in the $\Sigma$-tail index $T^{S}$; so
\[
A\text{'s positions } S\times T\;=\;(\text{content-input }S)\times(\text{constant-branch
slice } T\subseteq T^{S}\text{ of the }\Sigma\text{-tail index}).
\]
The worker is the constant/diagonal slice of the \emph{aggregation tail} --- a statement
about $\Sigma$, not about $F/\mathbf 1$. This sharpens, rather than rescues, the negative
verdict of clause~(2): entanglement selects constant branches; it does not supply the
bicomodule content.

\section{Grant placement}\label{sec:grant}

The weld now has a precise home in $\Poly$: it \emph{is} a data-migration prafunctor
$\Delta T\text{-}\Set\to\Delta S\text{-}\Set$, $X\mapsto(T\times X)^{S}$ --- ``read the whole
input state $S$; return, for each input, an output token in $T$ and a value in $X$,'' with
the branch-map choice aggregated by $\Sigma$. This places the \emph{state axis} of the
compositional-correctness program inside the same framed bicategory $\Cat$ that houses the
directed (retrofunctor) axis --- the positive deliverable the note hoped for. But the
identification ``worker $=$ PRA content'' is \emph{false}: state-entanglement ($A$) and
data-aggregation (the $\Sigma$-tail) are genuinely different collapses, and the
agent-orchestration weld semantics must keep them distinct.

\section{Scope and honesty}\label{sec:honesty}

\begin{itemize}
\item Lemma~\ref{thm:reg-lemma} is standard category theory; the contribution is
clauses~(2)/(3), the explicit $\bar F_{B}=(T\times V)^{S}$ with its factorization, and the
dissolution of Gap~5.2.
\item Theorem~\ref{thm:represented} uses the two single-sided represented-functor formulas
($F_{G}$, $F_{H}$); these are the SGF $\Set[c]$ / Prop.~5.7 dictionaries applied to the
indiscrete case, where $c\text{-}\Set\simeq\Set$, representables are terminal, and component
polynomials are $y^{S}$. Each step is elementary but was \emph{derived}, not re-brute-forced
from the bicomodule-composition equalizer; the axiom check (\S\ref{sec:clause1}) and the
cardinality check (\S\ref{sec:represented}) are the computational corroboration.
\item This note decides exactly the three clauses of Conjecture~5.1 plus Gap~5.2; no claim
is made beyond them. In particular clause~(2) is a \emph{refutation}, and the
partial bridge (clauses 1 and 3) is not inflated into a full worker$\,=\,$content
identification.
\end{itemize}

\end{document}
